\documentclass[12pt,a4paper]{report}

% =========================================================
% PACKAGES
% =========================================================

\usepackage[
    a4paper,
    left=1in,
    right=1in,
    top=1in,
    bottom=1in
]{geometry}

\usepackage[T1]{fontenc}
\usepackage{lmodern}
\usepackage{graphicx}
\usepackage{tikz}

\usetikzlibrary{
    positioning,
    arrows.meta,
    shapes.geometric
}

\usepackage{titlesec}
\usepackage{fancyhdr}
\usepackage{hyperref}
\usepackage{amsmath}
\usepackage{amssymb}
\usepackage{float}
\usepackage{caption}
\usepackage{booktabs}
\usepackage{array}
\usepackage{tabularx}
\usepackage{enumitem}
\usepackage{listings}
\usepackage{xcolor}
\usepackage{setspace}
\usepackage{url}

% =========================================================
% GENERAL FORMATTING
% =========================================================

\linespread{1.3}

\setlength{\parindent}{0.5in}
\setlength{\parskip}{2pt}

\setlist{
    nosep,
    leftmargin=*
}

\captionsetup{
    font=small,
    labelfont=bf
}

% =========================================================
% CHAPTER / SECTION FORMATTING
% =========================================================

\titleformat{\chapter}[display]
{\normalfont\bfseries\centering}
{\Large CHAPTER \thechapter}
{8pt}
{\Large}

\titlespacing*{\chapter}
{0pt}
{-15pt}
{18pt}

\titleformat{\section}
{\large\bfseries}
{\thesection}
{1em}
{}

\titlespacing*{\section}
{0pt}
{10pt}
{5pt}

\titleformat{\subsection}
{\normalsize\bfseries}
{\thesubsection}
{1em}
{}

\titlespacing*{\subsection}
{0pt}
{7pt}
{3pt}

% =========================================================
% HYPERLINK SETTINGS
% =========================================================

\hypersetup{
    colorlinks=true,
    linkcolor=black,
    urlcolor=blue,
    citecolor=black,
    pdftitle={
        Student Placement Prediction System Using Machine Learning
    },
    pdfauthor={
        Pallavikrishna B
    }
}

% =========================================================
% HEADER AND FOOTER
% =========================================================

\pagestyle{fancy}
\fancyhf{}

\fancyfoot[C]{\thepage}

\renewcommand{\headrulewidth}{0pt}
\renewcommand{\footrulewidth}{0pt}

% =========================================================
% TABLE COLUMN
% =========================================================

\newcolumntype{Y}{
    >{\raggedright\arraybackslash}X
}

% =========================================================
% CODE LISTING SETTINGS
% =========================================================

\lstset{
    language=Python,
    basicstyle=\ttfamily\footnotesize,
    breaklines=true,
    frame=single,
    numbers=left,
    numberstyle=\tiny,
    stepnumber=1,
    numbersep=6pt,
    showstringspaces=false,
    tabsize=4,
    captionpos=b
}

% =========================================================
% IMAGE COMMAND
% All images are in the same folder as this .tex file.
% =========================================================

\newcommand{\projectfigure}[3]{%
\begin{figure}[!htbp]
    \centering
    \IfFileExists{#1}{%
        \includegraphics[
            width=#2\textwidth,
            height=0.40\textheight,
            keepaspectratio
        ]{#1}%
    }{%
        \fbox{%
            \parbox[c][3.8cm][c]{0.78\textwidth}{%
                \centering
                Figure placeholder\\[3pt]
                \texttt{#1}
            }
        }%
    }
    \caption{#3}
\end{figure}
}

% =========================================================
% DOCUMENT
% =========================================================

\begin{document}

% =========================================================
% PRELIMINARY PAGES
% No visible page numbers
% =========================================================

\pagenumbering{roman}
\pagestyle{empty}

% =========================================================
% TITLE PAGE
% =========================================================

\begin{titlepage}

\centering

% ---------------------------------------------------------
% College Name
% ---------------------------------------------------------

\vspace*{0.3cm}

{\large\textbf{RAJIV GANDHI INSTITUTE OF TECHNOLOGY}}

\vspace{0.15cm}

{\large\textbf{GOVERNMENT ENGINEERING COLLEGE}}

\vspace{0.15cm}

{\large\textbf{KOTTAYAM - 686 501}}

\vspace{0.35cm}

% ---------------------------------------------------------
% College Logo
% ---------------------------------------------------------

\IfFileExists{image.png}{
    \includegraphics[
        width=0.20\textwidth
    ]{image.png}
}

\vspace{0.45cm}

% ---------------------------------------------------------
% Department
% ---------------------------------------------------------

{\normalsize\textbf{DEPARTMENT OF COMPUTER APPLICATIONS}}

\vspace{0.35cm}

{\small 20MCA245 - MINI PROJECT}

\vspace{0.45cm}

% ---------------------------------------------------------
% Project Title
% ---------------------------------------------------------

{\normalsize\textbf{STUDENT PLACEMENT PREDICTION SYSTEM}}

\vspace{0.15cm}

{\small USING MACHINE LEARNING}

\vspace{0.65cm}

% ---------------------------------------------------------
% Submitted By
% ---------------------------------------------------------

{\small Submitted By}

\vspace{0.15cm}

{\normalsize\textbf{PALLAVIKRISHNA B}}

\vspace{0.15cm}

{\small\textbf{KTE25MCA2042}}

\vspace{0.7cm}

% ---------------------------------------------------------
% Guide
% ---------------------------------------------------------

{\small Under the Guidance of}

\vspace{0.15cm}

{\normalsize\textbf{Prof. JASSIM RAFEEK}}

\vspace{0.5cm}

% ---------------------------------------------------------
% KTU Logo
% ---------------------------------------------------------

\IfFileExists{ktu.png}{
    \includegraphics[
        width=0.16\textwidth
    ]{ktu.png}
}

\vfill

% ---------------------------------------------------------
% University
% ---------------------------------------------------------

{\small APJ ABDUL KALAM TECHNOLOGICAL UNIVERSITY}

\vspace{0.15cm}

{\small THIRUVANANTHAPURAM}

\vspace{0.3cm}

{\small October 2026}

\end{titlepage}

% =========================================================
% CERTIFICATE PAGE
% =========================================================

\newpage
\thispagestyle{empty}

\begin{center}

\vspace*{0.55cm}

% Department
{\normalsize\textbf{DEPARTMENT OF COMPUTER APPLICATIONS}}

\vspace{0.65cm}

% College
{\large\textbf{RAJIV GANDHI INSTITUTE OF TECHNOLOGY}}

\vspace{0.05cm}

{\large\textbf{KOTTAYAM}}

\vspace{0.55cm}

% RIT Logo
\IfFileExists{image.png}{
    \includegraphics[
        width=0.19\textwidth
    ]{image.png}
}

\vspace{1.05cm}

% Certificate heading
{\normalsize\textbf{CERTIFICATE}}

\end{center}

\vspace{0.35cm}

% ---------------------------------------------------------
% Certificate text
% ---------------------------------------------------------

\noindent
This is to certify that the project entitled
\textbf{``STUDENT PLACEMENT PREDICTION SYSTEM USING MACHINE
LEARNING''} is a bonafide work carried out by
\textbf{PALLAVIKRISHNA B} (Register No:
\textbf{KTE25MCA2042}) during the academic year
\textbf{2025 - 2026} in partial fulfillment of the requirements
for the award of the degree of
\textbf{Master of Computer Applications} of
\textbf{APJ Abdul Kalam Technological University,
Thiruvananthapuram, Kerala.}

\vspace{2.0cm}

% ---------------------------------------------------------
% Signature Section
% ---------------------------------------------------------

\vspace{1.5cm}

\noindent
\begin{tabular*}{\textwidth}{@{\extracolsep{\fill}}ll}

\begin{tabular}[t]{l}
\textbf{Prof. Jassim Rafeek}\\
(Project Guide)\\
Assistant Professor\\
Dept of Computer Applications
\end{tabular}
&
\begin{tabular}[t]{r}
\textbf{Dr. Vineetha S}\\
Head of Department\\
Dept of Computer Applications
\end{tabular}

\end{tabular*}

\newpage

% =========================================================
% DECLARATION
% =========================================================

\newpage
\thispagestyle{empty}

\begin{center}
    {\large\textbf{DECLARATION}}
\end{center}

\vspace{0.45cm}

\noindent
I, undersigned hereby declare that the project report entitled
\textbf{``STUDENT PLACEMENT PREDICTION SYSTEM USING MACHINE
LEARNING''}, submitted for partial fulfillment of the requirements
for the award of the degree of Master of Computer Applications of
the APJ Abdul Kalam Technological University, Kerala is a bonafide
work done by me under the supervision of my mentor,
\textbf{Prof. JASSIM RAFEEK}. This submission represents my ideas
in my own words and where ideas or words of others have been
included, I have adequately and accurately cited and referenced
the original sources. I also declare that I have adhered to the
ethics of academic honesty and integrity and have not
misrepresented or fabricated any data or idea or fact or source in
my submission. I understand that any violation of the above will
be a cause for disciplinary action by the institute and / or the
University and can also evoke penal action from the sources which
have thus not been properly cited or from whom proper permission
has not been obtained. This report has not been previously formed
as the basis for the award of any degree, diploma, or similar
title of any other University.

\vspace{3.8cm}

\noindent
Place: Kottayam
\hfill
\textbf{PALLAVIKRISHNA B}

\vspace{0.35cm}

\noindent
Date: 01/10/2026

\newpage

% =========================================================
% ACKNOWLEDGEMENT
% =========================================================

\newpage
\thispagestyle{empty}

\begin{center}
    {\large\textbf{ACKNOWLEDGEMENT}}
\end{center}

\vspace{0.45cm}

\noindent
I want to express my gratitude to everyone who has supported me
throughout the endeavour. First and foremost, I give thanks to
God Almighty for His mercy and blessings, for without His
unexpected direction, this would still be only a dream.

\vspace{0.15cm}

\noindent
I sincerely thank \textbf{Dr. Prince A}, Principal, Rajiv Gandhi
Institute of Technology, Kottayam, for providing the environment
in which this project could be completed.

\vspace{0.15cm}

\noindent
I express my sincere gratitude to \textbf{Dr. Vineetha S}, Head
of the Department of Computer Applications, for granting
permission and providing the necessary support for the successful
completion of this project.

\vspace{0.15cm}

\noindent
I am deeply grateful to my Project Guide,
\textbf{Prof. Jassim Rafeek}, Assistant Professor, Department of
Computer Applications, for his valuable guidance, continuous
support, constructive suggestions, and encouragement throughout
the project.

\vspace{0.15cm}

\noindent
My heartfelt thanks to my project coordinator,
\textbf{Dr. Reena Murali}, Associate Professor, Department of
Computer Applications, for her constructive suggestions and
inspiration throughout the project.

\vspace{0.15cm}

\noindent
Finally, I would like to take this chance to express my gratitude
to the faculty and technical staff of the Department of Computer
Applications.

\vspace{1.2cm}

\begin{flushright}
    \textbf{PALLAVIKRISHNA B}
\end{flushright}

% =========================================================
% ABSTRACT
% =========================================================

\newpage
\thispagestyle{empty}

\begin{center}
    {\large\textbf{ABSTRACT}}
\end{center}

\vspace{0.45cm}

\noindent
The placement process is an important part of a student's
academic journey. Traditional placement evaluation generally
depends on academic performance, skills, certifications,
communication abilities, technical knowledge, projects, and
other student-related factors. Manual evaluation of these
factors can be time-consuming and may not provide a consistent
way of estimating a student's placement readiness.

\vspace{0.15cm}

\noindent
The proposed \textbf{Student Placement Prediction System Using
Machine Learning} is a web-based application developed to
provide an automated prediction of a student's placement
readiness. The system collects academic information,
certifications, communication assessment results, technical
assessment results, project information, and other relevant
student details.

\vspace{0.15cm}

\noindent
A Random Forest classification algorithm is used as the machine
learning model. The system uses a dataset containing approximately
10,000 records for model development. The trained model achieved
an accuracy of approximately \textbf{91.9\%} during evaluation.
The trained model is stored as
\texttt{placement\_model\_old.pkl} and loaded by the Flask
application during prediction.

\vspace{0.15cm}

\noindent
In addition to machine learning prediction, the system uses
rule-based evaluation. The final prediction is calculated using
an \textbf{80\% machine learning score and 20\% rule-based score}.
Based on the resulting score, students are classified into four
categories: High, Medium, Developing, and Low.

\vspace{0.15cm}

\noindent
The system also includes student registration and login,
academic information management, certification verification,
communication assessment, technical assessment, company-related
assessment, prediction generation, and result display. The
application is developed using HTML, CSS, JavaScript, Python,
Flask, MySQL, and Scikit-learn.

\vspace{0.15cm}

\noindent
The proposed system provides an integrated platform for
collecting student information and generating an automated
placement prediction that can be used as a supportive tool for
students and administrators.

% =========================================================
% TABLE OF CONTENTS
% =========================================================

\newpage
\thispagestyle{empty}

\tableofcontents

% =========================================================
% LIST OF FIGURES
% =========================================================

\newpage
\thispagestyle{empty}

\listoffigures

% =========================================================
% LIST OF TABLES
% =========================================================

\newpage
\thispagestyle{empty}

\listoftables

% =========================================================
% START ARABIC PAGE NUMBERING FROM INTRODUCTION
% =========================================================

\clearpage

\pagenumbering{arabic}
\setcounter{page}{1}
\pagestyle{fancy}

% =========================================================
% CHAPTER 1
% =========================================================

\chapter{INTRODUCTION}

The placement process serves as a pivotal bridge transitioning students from
academic institutions into the professional software and technology industries.
Higher education institutions continually seek dependable mechanisms to assess
student readiness and provide targeted intervention before campus recruitment
drives commence. In modern recruitment, placement evaluation is multidimensional,
encompassing academic consistency, core technical capabilities, communicative
proficiency, and verified co-curricular achievements. Machine learning offers
powerful data-driven techniques to identify underlying patterns within
historical recruitment records and forecast the placement likelihood of current
students. This project presents the Student Placement Prediction System, a
comprehensive web-based platform that combines machine learning classification
with rule-based evaluation to deliver accurate, multi-category placement
readiness assessments and structured guidance for students and academic
mentors.

\section{Introduction}

Campus placement is an important activity in higher education
institutions. It provides students with opportunities to begin
their professional careers after completing their academic
programmes. Student placement readiness depends on several
factors including academic performance, technical knowledge,
communication skills, certifications, projects, and participation
in different assessments.

Traditional placement evaluation often involves collecting
student information manually and evaluating individual factors
separately. With the increasing availability of student data,
machine learning can be used to identify patterns in historical
placement data and generate predictions for new students.

The proposed project develops a web-based Student Placement
Prediction System that combines a machine learning model with
rule-based evaluation. The application provides a centralized
platform for collecting student information and generating
placement readiness predictions.


\section{Problem Statement}

Students possess different levels of academic performance,
technical knowledge, communication skills, certifications,
projects, and other attributes. A system is therefore required
to collect these factors and use them to provide an automated
placement prediction.

The problem addressed by this project is the development of a
web-based system capable of collecting relevant student data,
processing the data, applying machine learning and rule-based
logic, and presenting the resulting placement readiness
category.


\section{Scope}

The scope of the project includes:

\begin{itemize}
    \item student registration and authentication;
    \item collection of academic details;
    \item collection and verification of certifications;
    \item communication skill assessment;
    \item technical assessment;
    \item company-related information and assessment;
    \item machine learning based prediction;
    \item rule-based score calculation;
    \item hybrid prediction;
    \item prediction result display;
    \item administrative management.
\end{itemize}

\section{Relevance}

The system is relevant to students and educational institutions
because it provides a structured method for evaluating multiple
placement-related factors.

Students can use the system to understand how different factors
contribute to their placement readiness. Administrators can use
the system to manage student information and assessment results.

The system is intended as a supportive prediction tool and not
as a replacement for actual recruitment decisions.

\section{Objectives}

The main objectives of the project are:

\begin{enumerate}
    \item To develop a web-based student placement prediction
    system.

    \item To collect academic, certification, communication,
    technical, and project-related information.

    \item To train a Random Forest machine learning model using
    placement-related data.

    \item To implement rule-based evaluation for selected
    student attributes.

    \item To combine machine learning and rule-based scores
    using a hybrid prediction approach.

    \item To classify students into different placement readiness
    categories.

    \item To provide an easy-to-use interface for students and
    administrators.
\end{enumerate}


\section{Organization of the Report}

The report is organized into seven chapters.

Chapter 1 introduces the project, its background, motivation,
scope, objectives, and contributions.

Chapter 2 discusses the system study, existing system, literature
review, identified gap, proposed system, feasibility, and
requirements.

Chapter 3 presents the development methodology, Agile approach,
development stages, testing strategy, and risk management.

Chapter 4 explains the system design, architecture, activity
design, database design, system flowchart, modules, and data flow.

Chapter 5 describes the implementation of the frontend, backend,
database, machine learning model, assessments, certification
verification, and prediction modules.

Chapter 6 presents the results, testing, machine learning
analysis, hybrid prediction analysis, and system performance.

Chapter 7 concludes the project and discusses possible future
enhancements.

% =========================================================
% CHAPTER 2
% =========================================================

\chapter{SYSTEM STUDY AND REQUIREMENT ANALYSIS}

A thorough system study and requirement analysis form the foundation for
designing an effective educational assessment platform. This chapter
investigates the traditional procedures employed in campus placement
evaluations, highlighting their limitations in handling multi-faceted student
data and providing timely, automated feedback. To address these operational
constraints, a modern, automated Student Placement Prediction System is
proposed. The chapter further examines the technical, economic, and
operational feasibility of the proposed application, followed by an exhaustive
analysis of the functional, non-functional, software, and hardware requirements
essential for robust deployment.

\section{Existing System}

Traditional placement evaluation systems mainly depend on manual
collection and analysis of student information. Academic
performance may be considered separately from communication
skills, technical skills, certifications, and projects.

In many cases, students are required to maintain their
information in different documents or platforms. Placement
officers may also need to manually review and compare student
records.

Such approaches can make the evaluation process time-consuming
and may not provide an integrated prediction mechanism.

\section{Limitations of Existing System}

The major limitations include:

\begin{itemize}
    \item Manual collection of student information.
    \item Separate evaluation of different skill areas.
    \item Time-consuming analysis of student records.
    \item Limited use of historical placement data.
    \item Lack of an integrated prediction mechanism.
    \item Difficulty in maintaining large numbers of student
    records manually.
    \item Limited personalized feedback to students.
\end{itemize}


\section{Proposed System}

The proposed system is a web-based Student Placement Prediction
System that integrates a Flask application, MySQL database, and
Random Forest machine learning model.

The system follows the following general process:

\begin{enumerate}
    \item Student registration and login.
    \item Collection of student details.
    \item Collection of academic information.
    \item Certification submission and verification.
    \item Communication assessment.
    \item Technical assessment.
    \item Collection of relevant company assessment information.
    \item Machine learning prediction.
    \item Rule-based score calculation.
    \item Hybrid score generation.
    \item Final prediction classification.
    \item Result display.
\end{enumerate}

\section{Development Methodology}

The project follows an iterative development approach. The system
was developed by dividing the work into manageable modules such
as authentication, academic details, certifications,
assessments, prediction, and administration.

Each module was developed and tested before integration with the
complete system.

\section{Feasibility Study}

\subsection{Technical Feasibility}

The system uses commonly available technologies such as Python,
Flask, MySQL, HTML, CSS, JavaScript, and Scikit-learn. Therefore,
the system can be developed and executed on a standard computer.

\subsection{Economic Feasibility}

The major software technologies used in the project are freely
available. The system does not require specialized hardware for
basic development and testing.

\subsection{Operational Feasibility}

The application provides a web-based interface through which
students and administrators can interact with the system.
Therefore, the proposed system can be operated without requiring
specialized technical knowledge from normal users.

\section{Requirement Analysis}

\subsection{Functional Requirements}

The functional requirements include:

\begin{itemize}
    \item User registration.
    \item User login.
    \item Student profile management.
    \item Academic information entry.
    \item Certification submission.
    \item Certification verification.
    \item Communication assessment.
    \item Technical assessment.
    \item Company assessment.
    \item Machine learning prediction.
    \item Rule-based scoring.
    \item Hybrid prediction.
    \item Result display.
    \item Administrative management.
\end{itemize}

\subsection{Non-Functional Requirements}

The major non-functional requirements are:

\begin{itemize}
    \item Usability.
    \item Reliability.
    \item Maintainability.
    \item Security.
    \item Performance.
    \item Scalability.
\end{itemize}

\subsection{Software Requirements}

\begin{table}[H]
\centering
\caption{Software Requirements}

\begin{tabularx}{\textwidth}{p{0.30\textwidth}Y}
\toprule
\textbf{Software} & \textbf{Purpose} \\
\midrule
Python & Backend and machine learning development \\
Flask & Web application framework \\
MySQL & Database management \\
HTML & Web page structure \\
CSS & Web page styling \\
JavaScript & Client-side interaction \\
Scikit-learn & Machine learning implementation \\
Pandas & Dataset processing \\
NumPy & Numerical operations \\
Git & Version control \\
\bottomrule
\end{tabularx}
\end{table}

\subsection{Hardware Requirements}

The system can be developed using a standard personal computer
with:
\begin{itemize}
    \item Minimum 4 GB RAM.
    \item Dual-core or higher processor.
    \item Sufficient storage for development tools and datasets.
    \item Internet connection for downloading packages and
    deployment-related activities.
\end{itemize}

% =========================================================
% CHAPTER 3
% =========================================================

\chapter{DEVELOPMENT METHODOLOGY}

Developing an integrated web application that couples predictive machine
learning algorithms with interactive user interfaces necessitates a structured,
responsive software development methodology. This chapter outlines the
iterative Agile-Scrum framework adopted throughout the lifecycle of the
Student Placement Prediction System. By decomposing the overall application
into manageable product backlog items---ranging from user authentication and
academic records management to assessment engines and prediction
pipelines---the system was engineered through phased, incremental iterations.
The chapter also details the development stages, multi-level testing
strategies, and risk mitigation protocols employed to guarantee application
reliability and code quality.

\section{Introduction}

The project was developed using an iterative methodology. The
application was divided into different functional modules and
developed incrementally.

This approach made it possible to test individual components
before integrating them into the complete application.

\section{Agile and Scrum}

Agile development focuses on iterative development and continuous
improvement. The proposed project follows Agile principles by
dividing the development process into smaller tasks.

The major activities include:

\begin{itemize}
    \item Requirement identification.
    \item System design.
    \item Module development.
    \item Database development.
    \item Machine learning model development.
    \item Integration.
    \item Testing.
    \item Refinement.
\end{itemize}

\section{Product Backlog}

\begin{table}[H]
\centering
\caption{Major Product Backlog Items}

\begin{tabularx}{\textwidth}{p{0.12\textwidth}p{0.35\textwidth}Y}
\toprule
\textbf{No.} & \textbf{Task} & \textbf{Description} \\
\midrule
1 & Authentication & Registration and login functionality \\
2 & Student Details & Collection of personal information \\
3 & Academics & Collection of academic performance \\
4 & Certifications & Certification submission and verification \\
5 & Communication & Communication skill assessment \\
6 & Technical & Technical assessment \\
7 & Prediction & Machine learning and hybrid prediction \\
8 & Administration & Administrative management \\
9 & Testing & Functional and integration testing \\
\bottomrule
\end{tabularx}
\end{table}

\section{Development Stages}

The project was developed through the following stages:

\begin{enumerate}
    \item Requirement analysis.
    \item System architecture design.
    \item Database design.
    \item Frontend development.
    \item Backend development.
    \item Machine learning model development.
    \item Module integration.
    \item Testing.
    \item Final refinement.
\end{enumerate}

\section{Development Process}

The frontend was developed using HTML, CSS, and JavaScript.
The backend was implemented using Python and Flask.

MySQL was used for persistent storage of student and assessment
information.

The machine learning component was developed using Scikit-learn.
The Random Forest model was trained using historical placement
data and saved using Python pickle.

\section{Testing Strategy}

Testing was performed at different levels:

\begin{itemize}
    \item Unit testing.
    \item Functional testing.
    \item Database testing.
    \item Integration testing.
    \item User interface testing.
    \item Machine learning model testing.
\end{itemize}

\section{Risk Management}

Possible risks identified during development included:

\begin{itemize}
    \item Incorrect database connectivity.
    \item Invalid user input.
    \item Missing certification files.
    \item Machine learning model loading errors.
    \item Incorrect score calculation.
    \item Frontend-backend integration problems.
\end{itemize}

These risks were reduced through validation, modular
development, error handling, and repeated testing.

% =========================================================
% CHAPTER 4
% =========================================================

\chapter{SYSTEM DESIGN}

System design translates the functional and non-functional requirements into a
cohesive architectural and data model. This chapter elaborates on the structural
and behavioral design of the Student Placement Prediction System. A robust
three-tier architecture is established to delineate presentation, application
logic, and data storage layers. The chapter presents the detailed activity
flows, entity relationships, database schema, and system flowchart governing
user interactions. In addition, the modular decomposition of authentication,
profile management, assessment administration, and hybrid prediction calculation
is described, illustrating the complete data progression from input submission
to classification output.

\section{Introduction}

System design describes the architecture, modules, database,
workflow, and interactions of the proposed Student Placement
Prediction System.

\section{System Architecture}

The system follows a three-layer architecture consisting of the
presentation layer, application layer, and data/model layer.

\begin{itemize}
    \item \textbf{Presentation Layer:} HTML, CSS and JavaScript
    provide the user interface.

    \item \textbf{Application Layer:} Flask handles application
    logic, routing, authentication, assessment processing, and
    prediction.

    \item \textbf{Data and Model Layer:} MySQL stores application
    data and the trained Random Forest model performs machine
    learning prediction.
\end{itemize}

\projectfigure
{system_architecture_colored.png}
{0.80}
{System Architecture of the Student Placement Prediction System}

\section{Activity Design}

The main activities of the system include registration, login,
student details, academic information, certification verification,
communication assessment, technical assessment, aptitude/company
assessment, prediction, and result display.

\projectfigure
{activity_design_colored.png}
{0.80}
{Activity Design of the Student Placement Prediction System}

\section{Database Design}

The database is designed to store student information,
academic details, certifications, assessment results, and
prediction-related information.

MySQL is used as the relational database management system.

\projectfigure
{database_design.png}
{0.78}
{Database Design of the Student Placement Prediction System}

\section{System Flowchart}

The overall system flow begins with user registration or login.
After authentication, the student enters the required information,
completes the assessments, and submits the information for
prediction.

\projectfigure
{system_flowchart.png}
{0.78}
{System Flowchart of the Student Placement Prediction System}

\section{Module Design}

The major modules are:

\subsection{Authentication Module}

This module handles student registration and login. User
credentials are validated before access to the student dashboard
is provided.

\subsection{Student Information Module}

This module collects personal and academic information required
for the prediction system.

\subsection{Certification Module}

Students can submit up to five certifications. Certification
documents are uploaded for verification.

Each verified certification contributes 3 points, with a maximum
certification score of 15.

\subsection{Communication Assessment Module}

The communication assessment contains 10 multiple-choice
questions.

The questions are divided into:

\begin{itemize}
    \item Grammar.
    \item Vocabulary.
    \item Reading.
    \item Professional communication.
    \item Situational communication.
\end{itemize}

Each section contains two questions, with each question carrying
2 marks.

\subsection{Technical Assessment Module}

The technical assessment contains 10 multiple-choice questions
related to technical knowledge.

\subsection{Prediction Module}

The prediction module receives the required input features and
passes them to the trained Random Forest model.

A rule-based score is also calculated. The final score is
calculated using the hybrid prediction formula.

\section{Data Flow}

The general data flow is:

\begin{center}
\textbf{
Student Input
$\rightarrow$
Flask Backend
$\rightarrow$
Validation
$\rightarrow$
MySQL / ML Model
$\rightarrow$
Hybrid Prediction
$\rightarrow$
Result
}
\end{center}

\section{Prediction Categories}

The final prediction score is classified into four categories.

\begin{table}[H]
\centering
\caption{Prediction Categories}

\begin{tabular}{p{0.30\textwidth}
                p{0.30\textwidth}
                p{0.25\textwidth}}
\toprule
\textbf{Category} & \textbf{Score Range} & \textbf{Meaning} \\
\midrule
High & $\geq 75$ & Higher placement readiness score \\
Medium & $50$--$74.99$ & Moderate placement readiness score \\
Developing & $35$--$49.99$ & Developing placement readiness score \\
Low & $<35$ & Lower placement readiness score \\
\bottomrule
\end{tabular}
\end{table}

% =========================================================
% CHAPTER 5
% =========================================================

\chapter{IMPLEMENTATION}

The implementation phase brings together the architectural specifications,
software frameworks, database structures, and machine learning models into an
operational web platform. This chapter details the technical realisation of
the frontend interfaces using HTML5, CSS3, and JavaScript, alongside the backend
web services built using Python and the Flask framework. It describes the
integration with the MySQL relational database, the implementation of dynamic
communication and technical assessment modules, and the verification pipeline
for student certifications. Furthermore, the chapter covers the training,
serialization, loading, and evaluation of the Random Forest machine learning
classifier, as well as the implementation of the hybrid scoring algorithm that
combines predictive probabilities with rule-based assessment scores.

\section{Introduction}

The proposed system was implemented as a web application using
HTML, CSS, JavaScript, Python, Flask, MySQL, and Scikit-learn.

The implementation consists of frontend pages, backend routes,
database operations, assessment modules, certification
verification, machine learning prediction, and administrative
functionality.

\section{Frontend}

The frontend was developed using HTML, CSS, and JavaScript.

The interface provides pages for:

\begin{itemize}
    \item Registration.
    \item Login.
    \item Student details.
    \item Academic details.
    \item Dashboard.
    \item Certification verification.
    \item Communication assessment.
    \item Aptitude assessment.
    \item Company assessment.
    \item Technical assessment.
    \item Prediction.
    \item Result display.
\end{itemize}

The interface was designed to provide a consistent user
experience across the different modules.

\section{Frontend Components}

\begin{table}[H]
\centering
\caption{Frontend Components}

\begin{tabularx}{\textwidth}{p{0.32\textwidth}Y}
\toprule
\textbf{Technology / Component} & \textbf{Purpose} \\
\midrule
HTML & Structure of web pages \\
CSS & Styling and page layout \\
JavaScript & Client-side interaction and validation \\
Forms & Collection of student information \\
MCQ Interface & Communication and technical assessments \\
Result Interface & Display of prediction score and category \\
\bottomrule
\end{tabularx}
\end{table}

\section{Backend}

Python and Flask were used to implement the backend.

The backend is responsible for:

\begin{itemize}
    \item handling HTTP requests;
    \item validating input;
    \item managing sessions;
    \item communicating with MySQL;
    \item processing assessment results;
    \item loading the machine learning model;
    \item calculating hybrid prediction;
    \item rendering result pages.
\end{itemize}

\section{Backend Components}

\begin{table}[H]
\centering
\caption{Backend Components}

\begin{tabularx}{\textwidth}{p{0.35\textwidth}Y}
\toprule
\textbf{Component} & \textbf{Purpose} \\
\midrule
Flask & Web application framework \\
Python & Application and prediction logic \\
MySQL Connector & Database communication \\
Scikit-learn & Machine learning model \\
Pickle & Saving and loading trained model \\
Session Management & Maintaining logged-in user information \\
Validation & Checking submitted input \\
\bottomrule
\end{tabularx}
\end{table}

\section{Application Routes}

The major application routes include:

\begin{table}[H]
\centering
\caption{Major Application Routes}

\begin{tabularx}{\textwidth}{p{0.38\textwidth}Y}
\toprule
\textbf{Route / Module} & \textbf{Purpose} \\
\midrule
/register & Student registration \\
/login & Student authentication \\
/dashboard & Student dashboard \\
/academics & Academic details \\
/certifications & Certification information \\
/company\_selection & Company/domain selection \\
/communication & Communication assessment \\
/technical\_assessment & Technical assessment \\
/prediction & Final prediction \\
/admin & Administrative functions \\
\bottomrule
\end{tabularx}
\end{table}

\section{Authentication and Student Information}

The registration module allows a student to create an account.
The login module validates the credentials and establishes a
session.

After login, the student can access the dashboard and complete
the required information.

\section{Certification Verification}

The system allows students to submit up to five certifications.

Each certification contains information such as certification
name and uploaded document.

The administrator can verify or reject submitted certifications.

The available statuses are:

\begin{itemize}
    \item Pending.
    \item Verified.
    \item Rejected.
\end{itemize}

Each verified certification contributes 3 points and the total
certification score is capped at 15.

\section{Communication Assessment}

The communication assessment contains 10 MCQs.

The question distribution is:

\begin{table}[H]
\centering
\caption{Communication Assessment Structure}

\begin{tabular}{p{0.45\textwidth}
                p{0.25\textwidth}
                p{0.20\textwidth}}
\toprule
\textbf{Section} & \textbf{Questions} & \textbf{Marks} \\
\midrule
Grammar & 2 & 4 \\
Vocabulary & 2 & 4 \\
Reading & 2 & 4 \\
Professional & 2 & 4 \\
Situational & 2 & 4 \\
\midrule
\textbf{Total} & \textbf{10} & \textbf{20} \\
\bottomrule
\end{tabular}
\end{table}

Skipped questions receive zero marks.

\section{Technical Assessment}

The technical assessment contains 10 randomly selected
multiple-choice questions.

The result is converted into the required technical score used
by the prediction module.

\section{Dataset}

The machine learning model was developed using a dataset
containing approximately 10,000 student-related records.

The dataset contains features representing different academic,
skill, project, and placement-related attributes.

The dataset is divided into training and testing data.

A test size of 20\% is used in the model training process.

\section{Random Forest Model}

Random Forest is an ensemble machine learning algorithm that
combines multiple decision trees.

The model was implemented using
\texttt{RandomForestClassifier} from Scikit-learn.

The model uses:

\begin{itemize}
    \item Random Forest classification.
    \item \texttt{random\_state = 42}.
    \item 80\% training data.
    \item 20\% testing data.
\end{itemize}

The evaluated model achieved approximately \textbf{91.9\%}
accuracy.

\begin{equation}
Accuracy =
\frac{\text{Correct Predictions}}
{\text{Total Predictions}}
\times 100
\end{equation}

\section{Model Loading}

The trained model is saved in the file
\texttt{placement\_model\_old.pkl}.

The Flask application loads this model when prediction is
required.

\section{Feature Importance}

Feature importance is used to understand the relative contribution
of input features in the Random Forest model.

The major feature importance values observed in the project are:

\begin{table}[H]
\centering
\caption{Selected Feature Importance Values}

\begin{tabular}{p{0.55\textwidth}p{0.25\textwidth}}
\toprule
\textbf{Feature} & \textbf{Importance} \\
\midrule
Communication & 0.2418 \\
CGPA & 0.1732 \\
Previous Semester Performance & 0.1501 \\
Projects & 0.1297 \\
\bottomrule
\end{tabular}
\end{table}

These values represent the feature importance values obtained
from the trained Random Forest model.

\section{Rule-Based Prediction}

In addition to the machine learning model, the system applies
rule-based evaluation.

The rule-based component incorporates selected explicit
conditions and scores, including certification verification and
assessment-related factors.

\section{Hybrid Prediction}

The final prediction is calculated using a hybrid approach.

The machine learning component contributes 80\% and the
rule-based component contributes 20\%.

The formula is:

\[
P_{\text{final}} =
0.80P_{\text{ML}} +
0.20P_{\text{Rules}}
\]

where:

\begin{itemize}
    \item $P_{\text{final}}$ is the final prediction score.
    \item $P_{\text{ML}}$ is the machine learning score.
    \item $P_{\text{Rules}}$ is the rule-based score.
\end{itemize}

\section{Prediction Classification}

The final score is classified as follows:

\begin{itemize}
    \item High: score $\geq 75$.
    \item Medium: score $\geq 50$ and $<75$.
    \item Developing: score $\geq 35$ and $<50$.
    \item Low: score $<35$.
\end{itemize}

\section{Result Page}

The result page displays the final prediction score and category.

The page provides the student with a clear representation of the
prediction generated from the submitted information.

\section{Administrative Module}

The administrative module allows authorized users to manage
student-related information.

Administrative functions include:

\begin{itemize}
    \item viewing student information;
    \item checking submitted certifications;
    \item verifying certifications;
    \item rejecting certifications;
    \item reviewing assessment results;
    \item accessing prediction-related information.
\end{itemize}

% =========================================================
% CHAPTER 6
% =========================================================

\chapter{TESTING AND ANALYSIS}

Verification and performance analysis are critical to ascertain that the
integrated system behaves reliably under realistic conditions and produces
accurate predictive assessments. This chapter presents the systematic testing
methodologies applied across individual modules and the unified application,
covering user authentication, data persistence, assessment scoring, and end-to-end
workflow execution. Subsequently, a rigorous analytical evaluation of the
trained Random Forest classifier is presented, detailing model accuracy,
training-testing splits, and feature importance distributions. Finally, the
chapter evaluates the efficacy of the hybrid prediction approach, which
harmonises algorithmic pattern discovery with explicit educational evaluation
rules.

\section{Introduction}

The developed Student Placement Prediction System was tested
module by module and then as an integrated application.

The testing focused on authentication, database operations,
student information, certifications, assessments, prediction,
and result display.

\section{Complete System Workflow}

The complete system workflow is:

\begin{enumerate}
    \item Student registration.
    \item Student login.
    \item Student details entry.
    \item Academic details entry.
    \item Certification submission.
    \item Certification verification.
    \item Communication assessment.
    \item Technical assessment.
    \item Company-related assessment.
    \item Prediction processing.
    \item Hybrid score calculation.
    \item Prediction category generation.
    \item Result display.
\end{enumerate}

\section{Functional Testing}

Functional testing was performed to verify whether each major
system function worked according to its requirements.

\begin{table}[H]
\centering
\caption{Functional Testing}

\begin{tabularx}{\textwidth}
{p{0.27\textwidth}p{0.42\textwidth}Y}
\toprule
\textbf{Module} & \textbf{Test} & \textbf{Expected Result} \\
\midrule
Registration & Submit valid student details &
Account created successfully \\

Login & Enter valid credentials &
Student dashboard displayed \\

Academics & Submit academic details &
Details stored successfully \\

Certification & Upload certificate &
Certificate stored for verification \\

Communication & Complete MCQs &
Communication score calculated \\

Technical & Complete technical questions &
Technical score calculated \\

Prediction & Submit all required inputs &
Prediction generated \\

Result & Open result page &
Final score and category displayed \\
\bottomrule
\end{tabularx}
\end{table}

\section{Authentication Testing}

Authentication testing was carried out to verify registration,
login, validation, and session handling.

\begin{table}[H]
\centering
\caption{Authentication Testing}

\begin{tabularx}{\textwidth}{p{0.35\textwidth}Y}
\toprule
\textbf{Test Case} & \textbf{Result} \\
\midrule
Valid registration & Successful \\
Duplicate registration & Validation applied \\
Valid login & Successful \\
Invalid login & Error displayed \\
Logout & Session terminated \\
\bottomrule
\end{tabularx}
\end{table}

\section{Database Testing}

Database testing was performed to verify insertion, retrieval,
updating, and storage of student-related information.

The database was checked for:

\begin{itemize}
    \item student details;
    \item academic details;
    \item certification records;
    \item assessment scores;
    \item prediction-related information.
\end{itemize}

\section{Assessment Testing}

Communication and technical assessments were tested using
different combinations of correct, incorrect, and skipped
responses.

The system correctly assigns marks based on the submitted
responses. Skipped communication questions receive zero marks.

\section{Machine Learning Result}

The Random Forest model was trained using approximately 10,000
records.

The evaluated accuracy was approximately:

\[
\boxed{91.9\%}
\]

The model was saved as:

\begin{center}
\texttt{placement\_model\_old.pkl}
\end{center}

The saved model is loaded by the Flask backend during prediction.

\section{Feature Analysis}

The major feature importance values obtained from the model were:

\begin{table}[H]
\centering
\caption{Feature Analysis}

\begin{tabular}{p{0.55\textwidth}p{0.25\textwidth}}
\toprule
\textbf{Feature} & \textbf{Importance} \\
\midrule
Communication & 0.2418 \\
CGPA & 0.1732 \\
Previous Semester Performance & 0.1501 \\
Projects & 0.1297 \\
\bottomrule
\end{tabular}
\end{table}

The values indicate the relative contribution of these features
within the trained Random Forest model.

\section{Hybrid Prediction Analysis}

The hybrid prediction combines the machine learning score and
rule-based score.

\[
P_{\text{final}} =
0.80P_{\text{ML}} +
0.20P_{\text{Rules}}
\]

The use of a hybrid approach allows the system to incorporate
both patterns learned from the dataset and explicit application
rules.

\section{Prediction Categories}

The final prediction categories are:

\begin{table}[H]
\centering
\caption{Final Prediction Categories}

\begin{tabular}{p{0.30\textwidth}p{0.35\textwidth}}
\toprule
\textbf{Category} & \textbf{Final Score} \\
\midrule
High & $\geq 75$ \\
Medium & $50$--$74.99$ \\
Developing & $35$--$49.99$ \\
Low & $<35$ \\
\bottomrule
\end{tabular}
\end{table}

\section{System Performance}

The integrated system was tested by completing the complete
student workflow from registration to final prediction.

The frontend, backend, database, assessment, and prediction
components were integrated and tested together.

The system successfully provides:

\begin{itemize}
    \item student authentication;
    \item data collection;
    \item certification management;
    \item assessment processing;
    \item machine learning prediction;
    \item hybrid prediction;
    \item result generation.
\end{itemize}

\section{Advantages}

The proposed system provides:

\begin{itemize}
    \item Automated placement readiness prediction.
    \item Centralized student information.
    \item Integration of academic and skill-related factors.
    \item Certification verification and scoring.
    \item Machine learning based prediction.
    \item Rule-based evaluation.
    \item Hybrid prediction mechanism.
    \item Web-based accessibility.
    \item Structured result presentation.
\end{itemize}

\section{Limitations}

The current implementation has certain limitations:

\begin{itemize}
    \item Prediction quality depends on the quality and
    representativeness of the training dataset.

    \item The current model uses a fixed trained model.

    \item The prediction categories are based on predefined
    thresholds.

    \item The system does not currently provide company-specific
    prediction for every organization.

    \item Continuous automatic model retraining is not currently
    implemented.
\end{itemize}

% =========================================================
% CHAPTER 7
% =========================================================

\chapter{CONCLUSION AND FUTURE SCOPE}

The successful design, implementation, and evaluation of the Student Placement
Prediction System validate the effectiveness of integrating machine learning
models with interactive web applications in educational environments. This
chapter summarises the technical achievements of the project, including the
automated aggregation of student academic and assessment metrics, high-accuracy
Random Forest classification, and balanced hybrid prediction scoring. It reviews
the advantages offered by the system as an objective decision-support mechanism
for students and academic institutions, acknowledges current operational
limitations, and outlines prospective avenues for future enhancement such as
larger cross-institutional datasets, explainable AI, and cloud deployment.

\section{Conclusion}

The \textbf{Student Placement Prediction System Using Machine
Learning} was developed as a web-based application for automated
placement readiness prediction.

The system integrates student information, academic performance,
certifications, communication skills, technical skills, and
assessment results into a single platform.

A Random Forest classification model was developed using
approximately 10,000 records. The model achieved an evaluated
accuracy of approximately 91.9\%.

The system also implements a hybrid prediction mechanism in which
80\% of the final score is obtained from the machine learning
component and 20\% from rule-based evaluation.

The final score is classified into High, Medium, Developing, and
Low categories.

The project demonstrates the integration of machine learning with
a Flask web application and MySQL database. The resulting system
provides a structured platform for collecting student information,
conducting assessments, managing certifications, and generating
placement readiness predictions.

\section{Future Scope}

The system can be enhanced in several ways.

\subsection{Larger Dataset}

The model can be trained using a larger and more diverse dataset
containing placement records from multiple academic institutions
and recruitment cycles.

\subsection{Additional Machine Learning Algorithms}

Additional classification algorithms can be implemented and
evaluated, such as:

\begin{itemize}
    \item Logistic Regression.
    \item Decision Tree.
    \item Support Vector Machine.
    \item Gradient Boosting.
    \item XGBoost.
\end{itemize}

\subsection{Explainable Prediction}

Future versions can provide explanations for individual
predictions by showing how different features contributed to the
result.

\subsection{Personalized Recommendations}

The system can provide personalized suggestions based on areas
where a student may need further preparation.

\subsection{Company-Specific Prediction}

The prediction system can be extended to consider specific
company requirements and job roles.

\subsection{Cloud Deployment}

The application can be deployed on cloud platforms to make it
accessible through the Internet.

\subsection{Advanced Administrative Dashboard}

The administrative dashboard can be enhanced with charts,
statistics, filtering, and reporting features.

\subsection{Automated Model Retraining}

The system can be extended to automatically retrain the machine
learning model using newly collected placement data.

\section{Final Remarks}

The project demonstrates how machine learning, web technologies,
and database systems can be integrated to develop an educational
decision-support application.

The proposed system provides a foundation that can be further
extended with additional data, improved prediction techniques,
explainability, personalized recommendations, and cloud-based
deployment.

% =========================================================
% REFERENCES
% =========================================================

\clearpage

\renewcommand{\bibname}{References}

\addcontentsline{toc}{chapter}{References}

\begin{thebibliography}{99}

\bibitem{breiman}
L. Breiman,
``Random Forests,''
\textit{Machine Learning},
Vol. 45, pp. 5--32, 2001.

\url{https://doi.org/10.1023/A:1010933404324}

\bibitem{sklearn}
Scikit-learn Developers,
\textit{Scikit-learn: Machine Learning in Python}.

\url{https://scikit-learn.org/stable/}

\bibitem{randomforest}
Scikit-learn Documentation,
\textit{RandomForestClassifier Documentation}.

\url{https://scikit-learn.org/stable/modules/generated/sklearn.ensemble.RandomForestClassifier.html}

\bibitem{flask}
Pallets Projects,
\textit{Flask Documentation}.

\url{https://flask.palletsprojects.com/}

\bibitem{python}
Python Software Foundation,
\textit{Python Documentation}.

\url{https://docs.python.org/3/}

\bibitem{mysql}
Oracle,
\textit{MySQL 8.4 Reference Manual}.

\url{https://dev.mysql.com/doc/refman/8.4/en/}

\bibitem{javascript}
MDN Web Docs,
\textit{JavaScript Documentation}.

\url{https://developer.mozilla.org/en-US/docs/Web/JavaScript}

\bibitem{pandas}
Pandas Development Team,
\textit{Pandas Documentation}.

\url{https://pandas.pydata.org/docs/}

\bibitem{numpy}
NumPy Developers,
\textit{NumPy Documentation}.

\url{https://numpy.org/doc/}

\bibitem{mlbook}
T. Hastie, R. Tibshirani, and J. Friedman,
\textit{The Elements of Statistical Learning},
Springer.

\url{https://hastie.su.domains/ElemStatLearn/}

\bibitem{hands}
A. G\'{e}ron,
\textit{Hands-On Machine Learning with Scikit-Learn, Keras, and TensorFlow},
O'Reilly Media.

\end{thebibliography}

% =========================================================
% APPENDIX A
% =========================================================

\clearpage

\chapter*{APPENDIX A: GIT COMMIT SCREENSHOT}

\addcontentsline{toc}{section}{Appendix A: Git Commit Screenshot}

The following screenshot shows the Git commit and repository
information of the project.

\begin{figure}[!htbp]
    \centering
    \IfFileExists{git1.png}{
        \includegraphics[
            width=0.85\textwidth,
            height=0.48\textheight,
            keepaspectratio
        ]{git1.png}
    }{
        \fbox{
            \parbox[c][5cm][c]{0.80\textwidth}{
                \centering
                git1.png not found
            }
        }
    }
    \caption{Git Commit Screenshot}
\end{figure}

% =========================================================
% APPENDIX B
% =========================================================

\clearpage

\chapter*{APPENDIX B: SYSTEM SCREENSHOTS}

\addcontentsline{toc}{section}{Appendix B: System Screenshots}

\section*{Login Page}

\begin{figure}[!htbp]
    \centering
    \IfFileExists{login.png}{
        \includegraphics[
            width=0.68\textwidth,
            height=0.30\textheight,
            keepaspectratio
        ]{login.png}
    }{
        \fbox{
            \parbox[c][3.2cm][c]{0.70\textwidth}{
                \centering
                login.png not found
            }
        }
    }
    \caption{Login Page}
\end{figure}

\vspace{0.2cm}

\section*{Registration Page}

\begin{figure}[!htbp]
    \centering
    \IfFileExists{register.png}{
        \includegraphics[
            width=0.68\textwidth,
            height=0.30\textheight,
            keepaspectratio
        ]{register.png}
    }{
        \fbox{
            \parbox[c][3.2cm][c]{0.70\textwidth}{
                \centering
                register.png not found
            }
        }
    }
    \caption{Registration Page}
\end{figure}

\clearpage

\section*{Academic Details Page}

\begin{figure}[!htbp]
    \centering
    \IfFileExists{academics.png}{
        \includegraphics[
            width=0.68\textwidth,
            height=0.30\textheight,
            keepaspectratio
        ]{academics.png}
    }{
        \fbox{
            \parbox[c][3.2cm][c]{0.70\textwidth}{
                \centering
                academics.png not found
            }
        }
    }
    \caption{Academic Details Page}
\end{figure}

\vspace{0.2cm}

\section*{Dashboard}

\begin{figure}[!htbp]
    \centering
    \IfFileExists{dashboard.png}{
        \includegraphics[
            width=0.68\textwidth,
            height=0.30\textheight,
            keepaspectratio
        ]{dashboard.png}
    }{
        \fbox{
            \parbox[c][3.2cm][c]{0.70\textwidth}{
                \centering
                dashboard.png not found
            }
        }
    }
    \caption{Student Dashboard}
\end{figure}

\clearpage

\section*{Communication Assessment}

\begin{figure}[!htbp]
    \centering
    \IfFileExists{communication.png}{
        \includegraphics[
            width=0.68\textwidth,
            height=0.30\textheight,
            keepaspectratio
        ]{communication.png}
    }{
        \fbox{
            \parbox[c][3.2cm][c]{0.70\textwidth}{
                \centering
                communication.png not found
            }
        }
    }
    \caption{Communication Assessment Page}
\end{figure}

\vspace{0.2cm}

\section*{Aptitude Assessment}

\begin{figure}[!htbp]
    \centering
    \IfFileExists{aptitude.png}{
        \includegraphics[
            width=0.68\textwidth,
            height=0.30\textheight,
            keepaspectratio
        ]{aptitude.png}
    }{
        \fbox{
            \parbox[c][3.2cm][c]{0.70\textwidth}{
                \centering
                aptitude.png not found
            }
        }
    }
    \caption{Aptitude Assessment Page}
\end{figure}

\clearpage

\section*{Company Assessment}

\begin{figure}[!htbp]
    \centering
    \IfFileExists{company_assessment.png}{
        \includegraphics[
            width=0.68\textwidth,
            height=0.30\textheight,
            keepaspectratio
        ]{company_assessment.png}
    }{
        \fbox{
            \parbox[c][3.2cm][c]{0.70\textwidth}{
                \centering
                company_assessment.png not found
            }
        }
    }
    \caption{Company Assessment Page}
\end{figure}

\vspace{0.2cm}

\section*{Technical Assessment}

\begin{figure}[!htbp]
    \centering
    \IfFileExists{technical.png}{
        \includegraphics[
            width=0.68\textwidth,
            height=0.30\textheight,
            keepaspectratio
        ]{technical.png}
    }{
        \fbox{
            \parbox[c][3.2cm][c]{0.70\textwidth}{
                \centering
                technical.png not found
            }
        }
    }
    \caption{Technical Assessment Page}
\end{figure}

\clearpage

\section*{Prediction Page}

\begin{figure}[!htbp]
    \centering
    \IfFileExists{prediction.png}{
        \includegraphics[
            width=0.75\textwidth,
            height=0.38\textheight,
            keepaspectratio
        ]{prediction.png}
    }{
        \fbox{
            \parbox[c][4cm][c]{0.75\textwidth}{
                \centering
                prediction.png not found
            }
        }
    }
    \caption{Final Prediction Page}
\end{figure}

% =========================================================
% APPENDIX C
% =========================================================

\clearpage

\chapter*{APPENDIX C: SAMPLE CODE}

\addcontentsline{toc}{section}{Appendix C: Sample Code}

\section*{C.1 Random Forest Model Training}

\begin{lstlisting}[caption={Random Forest Model Training}]
import pandas as pd

from sklearn.model_selection import train_test_split
from sklearn.ensemble import RandomForestClassifier
from sklearn.metrics import accuracy_score
import pickle

data = pd.read_csv("placement_dataset.csv")

X = data.drop("Placement", axis=1)
y = data["Placement"]

X_train, X_test, y_train, y_test = train_test_split(
    X,
    y,
    test_size=0.20,
    random_state=42
)

model = RandomForestClassifier(
    random_state=42
)

model.fit(X_train, y_train)

y_pred = model.predict(X_test)

accuracy = accuracy_score(
    y_test,
    y_pred
)

print("Accuracy:", accuracy)

with open(
    "placement_model_old.pkl",
    "wb"
) as file:

    pickle.dump(model, file)

print("Model saved successfully.")
\end{lstlisting}

\vspace{0.3cm}

\section*{C.2 Loading the Trained Model}

\begin{lstlisting}[caption={Loading the Trained Model}]
import pickle

with open(
    "placement_model_old.pkl",
    "rb"
) as file:

    model = pickle.load(file)
\end{lstlisting}

\clearpage

\section*{C.3 Hybrid Prediction}

\begin{lstlisting}[caption={Hybrid Prediction Logic}]
ml_score = prediction_score

rule_score = rule_based_score

final_score = (
    0.80 * ml_score
    + 0.20 * rule_score
)

if final_score >= 75:

    prediction = "High"

elif final_score >= 50:

    prediction = "Medium"

elif final_score >= 35:

    prediction = "Developing"

else:

    prediction = "Low"

print("Final Score:", final_score)
print("Prediction:", prediction)
\end{lstlisting}

\vspace{0.3cm}

\section*{C.4 Flask Application}

\begin{lstlisting}[caption={Basic Flask Application Structure}]
from flask import Flask
from flask import render_template
from flask import request

app = Flask(__name__)

@app.route("/")
def home():

    return render_template(
        "index.html"
    )

@app.route(
    "/login",
    methods=["GET", "POST"]
)
def login():

    # Authentication logic

    return render_template(
        "login.html"
    )

@app.route(
    "/prediction",
    methods=["GET", "POST"]
)
def prediction():

    # Prediction processing

    return render_template(
        "prediction.html"
    )

if __name__ == "__main__":

    app.run(
        debug=True
    )
\end{lstlisting}

\clearpage

\section*{C.5 MySQL Database Connection}

\begin{lstlisting}[caption={MySQL Database Connection}]
import mysql.connector

connection = mysql.connector.connect(
    host="localhost",
    user="root",
    password="your_password",
    database="placement_prediction"
)

cursor = connection.cursor()

print(
    "Database connected successfully"
)
\end{lstlisting}

\vspace{0.3cm}

\section*{C.6 Certification Score Calculation}

\begin{lstlisting}[caption={Certification Score Calculation}]
verified_count = 0

for certificate in certificates:

    if certificate["status"] == "Verified":

        verified_count += 1

certification_score = (
    verified_count * 3
)

if certification_score > 15:

    certification_score = 15
\end{lstlisting}

\vspace{0.3cm}

\section*{C.7 Prediction Category Classification}

\begin{lstlisting}[caption={Prediction Category Classification}]
if final_score >= 75:

    category = "High"

elif final_score >= 50:

    category = "Medium"

elif final_score >= 35:

    category = "Developing"

else:

    category = "Low"
\end{lstlisting}

% =========================================================
% END OF DOCUMENT
% =========================================================

\end{document}
